% Pillar-5 (MIN-REPORT-RL) standalone abstract
\begin{abstract}
Papers comparing RL algorithms for LLM post-training report a \emph{label}---GRPO,
DAPO, Dr.GRPO, GSPO---and treat the surrounding software stack as an
implementation detail. Our position is that this convention is backwards: on the
evidence available to us, the stack is a material part of the reported result,
and a label without its stack is not a reproducible experimental
specification. Three exhibits motivate the reporting standard. First, a \emph{nominally} matched-configuration backend
audit---every labeled knob (group size, LoRA setup, learning rate, data, prompt
template, decoding, and seed) harmonized across frameworks, though as we disclose
the managed backend was \emph{not} matched at the base-checkpoint level---spanned
final training rewards from $5.0\%$ to $85.6\%$: a $17\times$ same-label swing
that cannot be assigned to the backend (the managed stack
also silently pinned a different base checkpoint, and that undisclosed bundling
is itself part of the exhibit). Second, the same label denotes different
experiments across stacks: ``DAPO'' run on an open trainer with true dynamic
sampling (a toy-scale arithmetic audit) produced a mean Zero-Variance Fraction
(ZVF) of $0.00$, while ``DAPO'' approximated on a closed stack (GSM8K) via an
asymmetric-clip surrogate produced mean ZVF $0.58$---same name, opposite
telemetry. Third, a 12-cell head-to-head of four labeled algorithms on a single
fixed stack landed within a $0.034$ band of one another (mean last-10 reward
$0.710$--$0.744$), a spread comparable to seed-level variation, suggesting
that label-only comparisons can hide stack-level variation. We therefore
propose an eight-item \emph{minimum reportable
stack}---loss form, reference-policy/KL handling, sampler/backend/precision
\textbf{(including base-checkpoint identity: revision/hash)}, per-step ZVF/GU
trajectory, group-size schedule, disjoint held-out split,
decontamination plus a parser-robustness probe, and held-out pass@$k$
curves alongside pass@1---each item earning its place via a
documented ranking-relevant lever. We call out base-checkpoint identity explicitly
because our backend audit showed an \emph{undisclosed} checkpoint swap is itself a
ranking-relevant lever: the ``stack'' effect is only interpretable once the base
checkpoint is pinned and reported alongside the framework. We formalize an RL-reproducibility threat
model with flip-risk grades R0--R5, specify a toolchain (manifest emitter, stack
diff with CI verdicts, registry, auditor badge, and citation-lever tracer) that
makes the standard enforceable, and close with a call to action for venues.
\paragraph{Absorbed satellite (retired).}
The former condensed community position (roster R06) is retired from the live
tree to \texttt{archive/absorbed/R06\_min\_report/}. Do not count R06 as a
separate venue paper.
\end{abstract}
\paragraph{August 2026 evidence boundary.} The standard requires provenance,
estimands, and missing-cell reporting; labels and selected checkpoints are not
substitutes for controlled comparative evidence.
